% =============================================================
% Bölüm 2: Bool Cebri ve Mantık Kapıları
% =============================================================
\chapter{Bool Cebri ve Mantık Kapıları}
\label{chap:bool-cebri}

\begin{tcolorbox}[
    colback=blokyesil!15,
    colframe=sinyalyesil!60!black,
    title={\textbf{Öğrenme Hedefleri}},
    fonttitle=\bfseries,
    rounded corners,
    boxrule=0.8pt
]
Bu bölümü tamamladığınızda aşağıdakileri yapabileceksiniz:
\begin{itemize}[nosep, leftmargin=*]
    \item İkili değişken ve mantık fonksiyonu kavramlarını tanımlamak.
    \item Bool cebrinin Huntington aksiyomlarını ve temel teoremlerini uygulamak.
    \item De Morgan teoremiyle Bool ifadelerini sadeleştirmek.
    \item Yedi temel mantık kapısının (AND, OR, NOT, NAND, NOR, XOR, XNOR) doğruluk
    tablolarını ve Bool ifadelerini yazmak.
    \item NAND/NOR kapılarının fonksiyonel tamlığının ASIC tasarımındaki önemini açıklamak.
\end{itemize}
\end{tcolorbox}

\section{Giriş}
\label{sec:bool-giris}

Önceki bölümde gördüğümüz ikili sayılar, bir bilgisayarın belleğinde \emph{durağan} veriyi
temsil eder: bir sayı, bir karakter, bir renk kodu. Ama bir işlemci sadece veri saklamaz ---
toplama yapar, karşılaştırır, dallanma kararı verir. Yani veriyi \emph{işleyen} bir mekanizmaya
ihtiyacımız var; bu mekanizma da yine $0$ ve $1$'lerle, ama bu kez statik değil, girişe bağlı
olarak çıkış üreten devrelerle kurulur. İşte bu devrelerin arkasındaki matematik, konumuz olan
\emph{Bool cebri}\index{Bool cebri}dir.

Sayısal donanım tasarımının bu ikinci temel taşı, 1854'te George Boole'un önerdiği ve 1938'de
Claude Shannon'ın anahtarlama devrelerine uyguladığı Bool cebri üzerine kuruludur. Bu bölümde
ikili mantığın aksiyomlarını, temel teoremlerini ve bunların fiziksel karşılığı olan mantık
kapılarını inceleyeceğiz. Buradaki her kavram, önce Logisim'de simüle edilecek, sonra Verilog
ile üç modelleme biçiminde yazılacak ve son olarak Icarus/Yosys akışıyla sentezlenecektir.

\section{İkili Değişkenler ve Mantık Fonksiyonları}
\label{sec:ikili-degiskenler}

\begin{tanim}{İkili Değişken:}
Yalnızca iki değer alabilen bir değişkene \emph{ikili değişken}\index{ikili değişken} denir.
Bu değerler genellikle $\{0, 1\}$, $\{\text{DOĞRU}, \text{YANLIŞ}\}$ ya da
$\{\text{YÜKSEK}, \text{DÜŞÜK}\}$ gerilim seviyeleri olarak yorumlanır. Sayısal devre
tasarımında $0$ ve $1$ gösterimi kullanılır.
\end{tanim}

\begin{tanim}{Mantık Fonksiyonu:}
$n$ adet ikili değişken $x_1, x_2, \ldots, x_n$ alan ve tek bir ikili çıktı üreten
$f: \{0,1\}^n \to \{0,1\}$ eşlemesine \emph{mantık fonksiyonu}\index{mantık fonksiyonu}
(Boole fonksiyonu) denir. Bir mantık fonksiyonu üç eşdeğer biçimde ifade edilebilir:
\textbf{doğruluk tablosu}, \textbf{Bool ifadesi} ve \textbf{lojik devre şeması}. İlerleyen
bölümlerde aynı fonksiyonu bu üç biçimde de üretip birbiriyle karşılaştıracağız.
\end{tanim}

\begin{bilginot}[Neden Üç Eşdeğer Gösterim?]
Doğruluk tablosu bir fonksiyonun \emph{davranışını} kesin biçimde tanımlar ama büyüdükçe
($2^n$ satır) okunması zorlaşır. Bool ifadesi kısa ve cebirsel işlemlere (sadeleştirme,
De Morgan) uygundur. Devre şeması ise fiziksel gerçeklemeye en yakın gösterimdir. HDL
tasarımında yazdığımız Verilog kodu da sonunda bu üçüncü forma, yani bir kapı ağına
(netlist) sentezlenir.
\end{bilginot}

\section{Bool Cebrinin Aksiyomları}
\label{sec:huntington}

Bool cebri, Öklid geometrisi gibi \emph{aksiyomatik} bir sistemdir: az sayıda temel varsayımdan
(aksiyom) yola çıkılır ve cebrin tüm teoremleri bu varsayımlardan mantıksal çıkarımla türetilir.
George Boole'un 1854'teki özgün formülasyonu, matematikçi E.~V.~Huntington tarafından 1904'te
altı \emph{postülata} indirgendi. Bu postülatlar, ikisi toplama ($+$) ikisi çarpma ($\cdot$)
için olmak üzere dört çiftten oluşur; bu eşleşme, ilerleyen kısımda göreceğimiz
\emph{düalite ilkesi}nin kaynağıdır.

\begin{tanim}{Huntington Postülatları:}
$B=\{0,1\}$ kümesi, $+$ (OR) ve $\cdot$ (AND) ikili işlemleri ile $'$ (NOT) tekli işlemi
altında, her $x,y,z \in B$ için aşağıdaki altı özdeşliği sağlıyorsa \emph{Bool cebri} adını
alır:
\end{tanim}

\begin{table}[h]
\centering
\caption{Huntington Postülatları}
\label{tab:huntington}
\begin{tabular}{clll}
\toprule
No & Postülat & Postülat (düali) & Adı \\
\midrule
1 & $x+0=x$ & $x\cdot 1=x$ & Birim eleman \\
2 & $x+x'=1$ & $x\cdot x'=0$ & Tümleme \\
3 & $x+y=y+x$ & $x\cdot y=y\cdot x$ & Değişme \\
4 & $x+(yz)=(x+y)(x+z)$ & $x(y+z)=xy+xz$ & Dağılma \\
\bottomrule
\end{tabular}
\end{table}

\begin{bilginot}[Anahtarlama Devresi Yorumu]
Shannon'ın 1938'de gösterdiği gibi, bu postülatlar fiziksel bir anahtarlama devresi olarak
okunabilir: $x+y$ iki anahtarın \textbf{paralel} bağlanması (biri kapalıysa akım geçer),
$x\cdot y$ \textbf{seri} bağlanması (ikisi de kapalı olmalı), $x'$ ise $x$'in tersidir (biri
kapalıyken diğeri açıktır).
\end{bilginot}

\begin{ornek}[Postülat 2'nin Anahtarlama Devresiyle Doğrulanması]
$x$ ve $x'$ tanım gereği her zaman zıt durumdadır: biri kapalıyken diğeri mutlaka açıktır.
\begin{itemize}[nosep, leftmargin=*]
    \item \textbf{Seri bağlantı ($x\cdot x'$):} Akım geçmesi için \emph{ikisinin de} kapalı
    olması gerekir. Zıt oldukları için bu hiçbir zaman gerçekleşmez $\Rightarrow x\cdot x'=0$.
    \item \textbf{Paralel bağlantı ($x+x'$):} Akım geçmesi için \emph{en az birinin} kapalı
    olması yeterlidir. Zıt oldukları için her zaman biri kapalıdır $\Rightarrow x+x'=1$.
\end{itemize}
\end{ornek}

\begin{alistirma}[\kolay]
Postülat 1'i ($x+0=x$ ve $x\cdot 1=x$) yukarıdaki örnekteki gibi anahtarlama devresi
yorumuyla açıklayınız: $0$ değerindeki bir anahtar her zaman açık, $1$ değerindeki bir anahtar
her zaman kapalıdır. Bu anahtarın devreye paralel ya da seri eklenmesi neden devrenin
davranışını değiştirmez?
\end{alistirma}

\noindent\textbf{Çözüm:} $0$ değerindeki bir anahtar her zaman \emph{açıktır}, yani hiçbir zaman
akım geçirmez. Böyle bir anahtar $x$'e \textbf{paralel} eklendiğinde ($x+0$), devreye ek bir
akım yolu sağlamaz. Akımın geçip geçmemesi hâlâ tamamen $x$'e bağlıdır, dolayısıyla
$x+0=x$. $1$ değerindeki bir anahtar ise her zaman \emph{kapalıdır}, yani her zaman akım
geçirir. Böyle bir anahtar $x$'e \textbf{seri} eklendiğinde ($x\cdot 1$), akımı hiçbir zaman
engellemez --- yine akımın geçip geçmemesi tamamen $x$'e bağlı kalır, dolayısıyla $x\cdot 1=x$.

\begin{bilginot}[Sıradan Cebirden Fark]
Postülat 4'ün ilk yarısı, $x+(yz)=(x+y)(x+z)$, sıradan cebirde \textbf{geçerli değildir}
(ör. $x{=}2,y{=}3,z{=}4$ için sol taraf $14$, sağ taraf $30$ verir). Bool cebrinde geçerli
olmasının nedeni, değişkenlerin yalnızca $\{0,1\}$ değerlerini alması: bu sonlu kümede tüm
$2^3=8$ kombinasyon tek tek denenerek özdeşlik doğrulanabilir. Bu ispat yöntemine
\emph{mükemmel tümevarım} (perfect induction) denir ve Bool cebrinde sıkça kullanılır.
\end{bilginot}

\section{Düalite İlkesi ve Temel Teoremler}
\label{sec:dualite}

\begin{tanim}{Düalite İlkesi:}
Bir Bool özdeşliğinin \emph{düali}, ifadedeki her $+$ ile $\cdot$'nin, her $0$ ile $1$'in yer
değiştirilmesiyle elde edilir. Postülatların her biri kendi düaliyle çift geldiği için (bkz.
Tablo~\ref{tab:huntington}), bu çiftlenme \textbf{tüm türetilmiş teoremlere de yansır}: bir teoremi
postülatlardan kanıtladığınızda, düalini kanıtlamak için hiçbir ek işlem gerekmez. Kanıttaki
her satırda aynı yer değiştirmeyi uygulamak yeterlidir.
\end{tanim}

Bunu somutlaştırmak için, listedeki ilk teoremi doğrudan postülatlardan türetelim.

\begin{ornek}[Teorem 1'in Postülatlardan Türetilmesi: $x+x=x$]
\begin{align*}
x+x &= (x+x)\cdot 1 & &\text{[Postülat 1: } y\cdot 1=y \text{, burada } y=x+x \text{]}\\
    &= (x+x)(x+x') & &\text{[Postülat 2: } 1=x+x' \text{]}\\
    &= x+x\cdot x' & &\text{[Postülat 4: } (x+y)(x+z)=x+yz \text{]}\\
    &= x+0 & &\text{[Postülat 2: } x\cdot x'=0 \text{]}\\
    &= x & &\text{[Postülat 1: } x+0=x \text{]}
\end{align*}
\end{ornek}

\begin{bilginot}[Düalitenin Getirisi: İkinci Teorem Bedava]
Yukarıdaki kanıtın her satırında $+ \leftrightarrow \cdot$ ve $0 \leftrightarrow 1$ değişimini
uygularsanız, $x\cdot x=x$ özdeşliğinin kanıtını elde edersiniz --- tek bir ek adım atmadan.
Bu yüzden Bool cebri kaynaklarında teoremler genellikle çift çift, tek bir kanıtla sunulur.
\end{bilginot}

Aynı türetme yöntemiyle elde edilen diğer temel teoremler, ileride sadeleştirme
işlemlerinde sürekli kullanacağımız için aşağıda topluca listelenmiştir.

\begin{table}[h]
\centering
\caption{Temel Teoremler}
\begin{tabular}{clll}
\toprule
No & Teorem & Teorem (düali) & Adı \\
\midrule
1 & $x+x=x$ & $x\cdot x=x$ & Yinelenirlik (idempotence) \\
2 & $x+1=1$ & $x\cdot 0=0$ & Yutan eleman \\
3 & $(x')'=x$ & --- & Çifte tümleme \\
4 & $x+xy=x$ & $x(x+y)=x$ & Yutma (absorption) \\
5 & $x+(y+z)=(x+y)+z$ & $x(yz)=(xy)z$ & Birleşme (associative) \\
\bottomrule
\end{tabular}
\end{table}

\begin{alistirma}[\orta]
Teorem 4'ü ($x+xy=x$) yukarıdaki Teorem 1 kanıtını örnek alarak postülatlardan türetiniz.
\emph{İpucu:} $x$'i önce $x\cdot 1$ olarak yazıp Postülat 4'ü (dağılma) tersten uygulamayı
deneyin.
\end{alistirma}

\noindent\textbf{Çözüm:}
\begin{align*}
x+xy &= x\cdot 1+x\cdot y & &\text{[Postülat 1: } x=x\cdot 1 \text{]}\\
     &= x\cdot(1+y) & &\text{[Postülat 4: } xy_1+xy_2=x(y_1+y_2) \text{, dağılma]}\\
     &= x\cdot 1 & &\text{[Teorem 2: } 1+y=1 \text{]}\\
     &= x & &\text{[Postülat 1: } x\cdot 1=x \text{]}
\end{align*}
